\section{Introdução}

O avanço das redes de computadores e a crescente digitalização de serviços têm
ampliado significativamente a exposição de sistemas a ameaças cibernéticas. Esse
crescimento é acompanhado por um aumento expressivo no volume de dados trafegados,
que passou de aproximadamente 16 GB por usuário ao mês em 2017 para cerca de 50 GB
mensais em 2022 \cite{Thakkar2021Survey}. Esse cenário amplia a superfície de
ataque e se reflete no aumento da quantidade e da sofisticação de ameaças, como
negações de serviço, varreduras de rede e acessos não autorizados. Paralelamente, os
impactos de falhas de segurança tornam-se cada vez mais críticos: mais de
14 bilhões de registros de dados foram vazados desde 2013, além de prejuízos
financeiros como os 29,8 bilhões de dólares perdidos em golpes
telefônicos apenas no ano de 2020 \cite{zehedi2023comparative}.

Os Sistemas de Detecção de Intrusão (IDS) monitoram atividades suspeitas e
identificam possíveis ataques por meio da análise de tráfego e de comportamento.
O conceito surgiu na década de 1980, baseado na análise de \en{logs} e de padrões
conhecidos, e evoluiu na década de 1990 para técnicas de detecção por anomalia,
capazes de identificar ataques desconhecidos \cite{patrick2022study}. Desde então,
os IDS consolidaram-se como mecanismos essenciais para a proteção de ativos
digitais \cite{emad2022intrusion}.

Nesse sentido, o uso de técnicas de Aprendizado de Máquina (\en{Machine Learning},
ML) tem se consolidado como alternativa promissora \cite{Khraisat2019Survey}.
Esses métodos identificam padrões complexos nos dados sem depender exclusivamente
de regras pré-definidas. A literatura, contudo, é dominada por avaliações em lote
sobre bases rotuladas, nas quais acurácias acima de 99\% são comuns. Levar um
modelo desses para uma rede real expõe problemas que a avaliação em lote não
mede: não há rótulo no tráfego de produção, o tráfego legítimo muda com o tempo,
o erro mais caro é o falso negativo silencioso, e o próprio processo de retreino
pode ser explorado por um atacante paciente \cite{Sommer2010Outside, Biggio2012Poisoning}.

Este trabalho tem como objetivo geral projetar, implementar e avaliar um Sistema de
Detecção de Intrusão em Redes (NIDS) baseado em Aprendizado de Máquina, evoluindo
de um estudo exploratório experimental em lote (MVP1) para uma arquitetura
operacional conteinerizada em tempo real, apta a monitorar o tráfego contínuo e
dotada de mecanismos de defesa contra o envenenamento de modelos (MVP2). Para
alcançar o objetivo geral, definem-se os seguintes objetivos específicos:

\begin{itemize}
    \item \textbf{Estudo exploratório (MVP1):} comparar algoritmos supervisionados
    sobre a base NSL-KDD e medir o efeito da redução de dimensionalidade (correlação,
    PCA e seleção por importância) sobre desempenho e custo computacional.
    \item \textbf{Classificação em dois estágios encadeados (MVP2):} um Estágio 1
    não supervisionado (\en{Isolation Forest}) para filtrar anomalias em tempo real,
    acompanhado de uma árvore de decisão \en{surrogate} que gera regras legíveis para
    cada alerta, e um Estágio 2 supervisionado (SVM linear) que refina os alertas a
    partir dos vereditos do analista, em regime de aprendizado ativo.
    \item \textbf{Resiliência ao envenenamento:} sustentar o retreino em um
    \en{pool} curado, com quarentena de 7 dias, peso de recência com meia-vida de
    14 dias, despejo de janelas comprometidas e rastreabilidade de quais modelos
    absorveram cada janela.
    \item \textbf{Portão de promoção:} avaliar todo candidato contra um conjunto
    rotulado de referência (\en{golden set}), exigindo a preservação do \en{recall}
    por cenário de ataque e limitando a taxa de falso positivo no controle benigno
    a 2\% antes de qualquer promoção para produção.
    \item \textbf{Infraestrutura de monitoramento:} empacotar o sistema em
    contêineres \en{Docker} sem privilégio de \en{root}, com captura passiva do
    tráfego espelhado, persistência em PostgreSQL/TimescaleDB e painéis em Grafana.
\end{itemize}

Este trabalho está estruturado em cinco seções principais. Além desta Introdução,
a Seção II apresenta a Fundamentação Teórica e os trabalhos relacionados. A Seção III
detalha a Metodologia, contemplando o estudo em lote (MVP1) e a arquitetura
operacional (MVP2). A Seção IV apresenta e discute os resultados das duas fases.
Por fim, a Seção V sintetiza as conclusões e indica trabalhos futuros.
